\section{Population Model and Induced Network}
\label{sec:model}

\subsection{Assignment state (relaxed robot counts $z$)}
A team of $N=\sum_\tau n_\tau$ robots is partitioned into $T$ types
$\tau\in\{1,\dots,T\}$ (heterogeneous in force capacity and communication range),
with $n_\tau$ robots of type $\tau$. There are $M$ loads with $H$ contact slots
each, a set $\mathcal R$ of candidate relay sites, and the common action set
\begin{equation}
\mathcal A=\underbrace{\{(k,h)\}_{k\le M,\,h\le H}}_{\text{load--slot}}
\;\cup\;\mathcal R\;\cup\;\{\mathrm{o}\},
\label{eq:actions}
\end{equation}
with $\mathrm o$ idle. We write $a\in\mathcal A$ for a generic action, $(k,h)$ for a
load--slot, and $r\in\mathcal R$ for a relay site.

The \emph{assignment state} is the relaxed-count matrix
$z=(z_{\tau a})\in\R_{\ge0}^{T\times|\mathcal A|}$: $z_{\tau a}$ is a continuous
relaxation of the number of type-$\tau$ robots assigned to action $a$---neither a
probability nor a physical force. Each type spends its budget, and each slot or
relay site holds at most one robot,
\begin{equation}
\begin{aligned}
&\textstyle\sum_{a\in\mathcal A} z_{\tau a}=n_\tau\ \ \forall\tau,
\qquad z_{\tau a}\ge0,\\[1mm]
&\textstyle\sum_\tau z_{\tau,kh}\le1\ \ \forall(k,h),
\qquad \sum_\tau z_{\tau,r}\le1\ \ \forall r.
\end{aligned}
\label{eq:conservation}
\end{equation}
The feasible set
$\mathcal X=\{z\in\R_{\ge0}^{T\times|\mathcal A|}:\eqref{eq:conservation}\}$ is a
compact convex polytope. A finite team instantiates $z$ by an integer assignment
(one action per robot); integerization is treated in Remark~\ref{rem:integer}. The
state $z$ is distinct from the physical poses $q\in(\mathrm{SE}(2))^N$, which enter
only the closed-loop simulation (Section~\ref{sec:exp}). Multiplying a count
$z_{\tau kh}$ by a per-robot force yields a total force, so every aggregate below is
dimensionally a force or torque.

\begin{assumption}[Standing hypotheses]\label{as:standing}
Throughout: \emph{(i)} $\mathcal X$ is compact and convex; \emph{(ii)} the aggregate
map $z\mapsto s(z)$ of \eqref{eq:tensor} is linear, so the wrench-feasible set
$\mathcal X_{\mathrm f}=\{z\in\mathcal X:s_{k\ell}(z)\ge d_{k\ell}\}$ is a polytope;
\emph{(iii)} the candidate-site Laplacian $\Lap(z)$ of \eqref{eq:laplacian} is affine
in $z$ with each elementary $L_e\succeq0$ and each coefficient $\ge0$; \emph{(iv)}
the communication graph is undirected and weighted, so every $\Lap(z)$ is a symmetric
graph Laplacian; and \emph{(v)} $\mathcal X_{\mathrm f}\neq\varnothing$.
\end{assumption}

\subsection{Support-function wrench certificate}
\emph{Physical force set.} A type-$\tau$ robot with heading $\phi$ exerts contact
forces in the heading-aligned ellipse $U_\tau(\phi)$ with longitudinal semi-axis
$F_\tau$ (along $\phi$) and lateral semi-axis $\kappa F_\tau$ ($\kappa\le1$, the
nonholonomic penalty). \emph{Certificate set.} The strategic layer does not track
headings, so for the certificate we use the heading-independent, centrally
symmetric \emph{inner} regular $2m$-gon $U_\tau^{\mathrm{in}}$ inscribed in the
disk of radius $\kappa F_\tau$. Since $\mathrm{disk}(\kappa F_\tau)\subseteq
U_\tau(\phi)$ for every $\phi$ (the ellipse's minimum semi-axis is $\kappa
F_\tau$), we have $U_\tau^{\mathrm{in}}\subseteq U_\tau(\phi)$ for all $\phi$;
feasibility tested against $U_\tau^{\mathrm{in}}$ is therefore \emph{conservative}
and cannot create a false positive.

Let $A_{kh}\in\R^{3\times2}$ be the contact-to-wrench map at slot $h$ of load $k$,
$\eta_{k\ell}\in\R^3$ ($\ell\le P$) a probing direction, and $h_{U}$ the support
function of $U$. Define the \emph{support-capability map}
$\Wmap\in\R^{T\times M\times H\times P}$ by
\begin{equation}
\Wmap_{\tau,kh\ell}=h_{U_\tau^{\mathrm{in}}}\!\big(A_{kh}^{\top}\eta_{k\ell}\big),
\qquad
d_{k\ell}=\eta_{k\ell}^{\top}\wdem_k ,
\label{eq:tensor}
\end{equation}
with $\wdem_k\in\R^3$ the demanded wrench of load $k$, and the population
directional capacity $s_{k\ell}(z)=\sum_{\tau,h}\Wmap_{\tau,kh\ell}\,z_{\tau kh}$
($z_{\tau kh}$ the count of type $\tau$ on slot $(k,h)$). Let
$\Wmap(C_k)=\bigoplus_{\tau,h}z_{\tau kh}A_{kh}U_\tau^{\mathrm{in}}$ be the
reachable population wrench set.

\begin{proposition}[Support-function wrench certificate]\label{prop:membership}
\emph{(i)} For every $\eta$,
$h_{\Wmap(C_k)}(\eta)=\sum_{\tau,h}z_{\tau kh}\,h_{U_\tau^{\mathrm{in}}}(A_{kh}^{\top}\eta)$.
\emph{(ii)} For any finite direction set $\mathcal D$,
$\{\eta^{\top}\wdem_k\le h_{\Wmap(C_k)}(\eta):\eta\in\mathcal D\}$ is
\emph{necessary} for $\wdem_k\in\Wmap(C_k)$. \emph{(iii)} it is \emph{sufficient}
iff $\mathcal D$ contains a complete set of outward facet normals of
$\Wmap(C_k)$, or an equivalent exact H-representation is used. \emph{(iv)} If every
$U_\tau$ is a polytope, $\Wmap(C_k)$ is a polytope. \emph{(v)} If every $U_\tau$ is
a zonotope, $\Wmap(C_k)$ is a zonotope. In particular each $U_\tau^{\mathrm{in}}$ is
a zonotope, so $\Wmap(C_k)$ is a zonotope, a complete finite $\mathcal D$ exists,
and exact membership is a linear program.
\end{proposition}
\begin{proof}
\emph{(i)} support functions are additive over Minkowski sums and satisfy the
linear-image identity $h_{MU}(\eta)=h_U(M^{\top}\eta)$. \emph{(ii)} supporting
hyperplane. \emph{(iii)} a compact convex body equals the intersection of its
supporting half-spaces at its outward facet normals; the finite subset $\mathcal D$
is exact precisely when it contains them. \emph{(iv)--(v)} linear images and
Minkowski sums of polytopes are polytopes, and of zonotopes are zonotopes; a
regular $2m$-gon is a zonotope. A general polytope is \emph{not} a zonotope (e.g.\
a triangle is not centrally symmetric), so we never infer zonotope from polytope.
\end{proof}
We call load $k$ \emph{wrench-feasible} at $z$ when $s_{k\ell}(z)\ge d_{k\ell}$ for
all $\ell$ (Assumption~\ref{as:standing}(ii)); by
$U_\tau^{\mathrm{in}}\subseteq U_\tau(\phi)$ this certificate is conservative---a
certified wrench is realizable by the physical elliptical forces. Unlike a scalar
capacity, \eqref{eq:tensor} separates force limit, lever arm, and slot.

\subsection{Decision-induced candidate-site Laplacian $\Lap(z)$}
We separate two graphs throughout. The \emph{decision-induced candidate-site
Laplacian} $\Lap(z)$ depends only on the strategic state $z$: a fixed family of
candidate communication links $e$ over $V$ nodes is switched on by the roles the
composition assigns. Only robots that do not lift can relay; with
$\mathcal C_{\tau re}\ge0$ the contribution of a type-$\tau$ mass in role $r$ to
link $e$,
\begin{equation}
a_e(z)=\sum_{\tau,r}\mathcal C_{\tau re}\,z_{\tau r},
\qquad
\Lap(z)=\sum_e a_e(z)\,L_e\in\R^{V\times V},
\label{eq:laplacian}
\end{equation}
which is affine in $z$ (Assumption~\ref{as:standing}(iii)). This is a convex
\emph{surrogate}: the actual, \emph{geometric} Laplacian $\Lgeo(q)$, whose weights
depend on inter-robot distances through the physical poses $q$, is nonconvex in
$q$ and is \emph{not} covered by the convex theory of
Sections~\ref{sec:wise}--\ref{sec:algorithm}; it is exercised only in the
closed-loop simulation, where we report the empirical agreement between
$\Lgeo(q)$ and its surrogate $\Lap(z)$. All theorems below use $\Lap(z)$ only.

Let $Q\in\R^{V\times(V-1)}$ be an orthonormal basis of $\mathbf 1^{\perp}$; the
information constraint is the LMI
\begin{equation}
Q^{\top}\Lap(z)\,Q\succeq\sigmastar I
\;\Longleftrightarrow\;
\lamtwo(\Lap(z))\ge\sigmastar,
\label{eq:lmi}
\end{equation}
affine in $z$, with PSD matrix dual $Z\succeq0$; at a simple minimiser
$Z=\pi vv^{\top}$ recovers a scalar connectivity price $\pi$
(Section~\ref{sec:algorithm}). The threshold $\sigmastar$ is derived, not assumed,
in Theorem~\ref{thm:stability}.

\begin{lemma}[Convex information set]\label{lem:convex}
$z\mapsto\lamtwo(\Lap(z))$ is concave, so \eqref{eq:lmi} defines a convex set.
\end{lemma}
\begin{proof}
$\lamtwo(L)=\min_{\|v\|=1,\,v\perp\mathbf 1}v^{\top}Lv$ is a pointwise minimum of
functions linear in $L$~\cite{fiedler1973}, hence concave; composed with the affine
$\Lap(z)$ it stays concave, and its super-level set is
convex~\cite{ghosh2006growing}.
\end{proof}

\begin{lemma}[Geometric transfer under model mismatch]\label{lem:bridge}
Suppose robots track their intended candidate sites so that
$\|\Lgeo(q)-\Lap(z)\|_2\le\varepsilon_L$. Then Weyl's inequality gives
\begin{equation}
\lamtwo(\Lgeo(q))\;\ge\;\lamtwo(\Lap(z))-\varepsilon_L .
\label{eq:weyl}
\end{equation}
Consequently, with the design margin
$\delta=\varepsilon_L+\varepsilon_{\mathrm{est}}+\delta_{\mathrm{num}}$
($\varepsilon_L$: geometric/model mismatch; $\varepsilon_{\mathrm{est}}$:
distributed eigenvalue-estimation error; $\delta_{\mathrm{num}}$: numerical
margin), the WISE requirement $\lamtwo(\Lap(z))\ge\sigmastar:=\sigmadyn+\delta$
implies $\lamtwo(\Lgeo(q))\ge\sigmadyn$, the physical stability threshold.
\end{lemma}
\begin{proof}
Weyl's inequality for symmetric matrices gives
$|\lamtwo(\Lgeo)-\lamtwo(\Lap)|\le\|\Lgeo-\Lap\|_2\le\varepsilon_L$, which is
\eqref{eq:weyl}; subtracting $\varepsilon_{\mathrm{est}}+\delta_{\mathrm{num}}$
absorbs the estimator and numerical errors.
\end{proof}
The convex theory of Sections~\ref{sec:wise}--\ref{sec:algorithm} is stated over
the affine $\Lap(z)$; the physical guarantee transfers to the geometric $\Lgeo(q)$
\emph{only} through \eqref{eq:weyl} under the margin $\delta$, and is otherwise
reported as simulation evidence. The margin's components are not chosen
arbitrarily: $\varepsilon_L$ is measured (Section~\ref{sec:exp}, bounded by the
Lipschitz constant $C_L\rho$ of the weight model), and $\varepsilon_{\mathrm{est}}$,
$\delta_{\mathrm{num}}$ are documented in the generated certificate.